\chapter{2019:John B. Goodenough and M. Stanley Whittingham and Akira Yoshino}

\label{chap:chemistry2019}

The Nobel Prize in Chemistry for the year 2019 was awarded jointly to John B. Goodenough and M. Stanley Whittingham and Akira Yoshino.

\textbf{Motivation:}

for the development of lithium-ion batteries

\section{John B. Goodenough}

\subsection*{Early Life and Education}

John B. Goodenough was born in 1922 in Jena, Germany.

\subsection*{Career and Major Research}

Goodenough earned a bachelor's degree in mathematics from Yale University in 1944 and a doctorate in physics from the University of Chicago in 1952. He worked at the Massachusetts Institute of Technology and the University of Oxford before becoming a professor at the University of Texas at Austin in 1986. In 1980 he identified lithium cobalt oxide (LiCoO2) as a cathode material that stores lithium ions at a higher voltage than earlier materials.

\subsection*{Nobel Prize-winning Work}

The work for which John B. Goodenough was awarded the Nobel Prize in Chemistry in 2019 was described by the Nobel Committee as: "for the development of lithium-ion batteries".

Goodenough's cobalt oxide cathode doubled the voltage of the battery and made a practical, high-energy rechargeable battery possible.

\subsection*{Significance and Impact}

His cathode was the key step that allowed lithium-ion batteries to store enough energy for portable electronics and electric vehicles. At the age of 97, Goodenough became the oldest recipient of a Nobel Prize.

\subsection*{Later Career and Honors}

Goodenough continued to work on battery materials at the University of Texas at Austin into his nineties. He died on 2023-06-25 in Austin, Texas, USA.

\subsection*{Legacy}

The rechargeable lithium-ion battery powers the world's mobile phones, laptops, and electric cars, and the layered-oxide cathodes he pioneered remain in widespread use.

\subsection*{Selected Publications}

"LixCoO2 (0<x$\leq$1): A New Cathode Material for Batteries of High Energy Density" (Materials Research Bulletin, 1980)

\clearpage

\section{M. Stanley Whittingham}

\subsection*{Early Life and Education}

M. Stanley Whittingham was born in 1941 in Nottingham, England.

\subsection*{Career and Major Research}

Whittingham received his doctorate from the University of Oxford in 1968. He worked at the Exxon research laboratories and later became a professor at the Binghamton University of the State University of New York. In the early 1970s he showed that lithium ions can move in and out of layered titanium disulfide, creating a rechargeable cathode with the intercalation mechanism used in lithium-ion batteries.

\subsection*{Nobel Prize-winning Work}

The work for which M. Stanley Whittingham was awarded the Nobel Prize in Chemistry in 2019 was described by the Nobel Committee as: "for the development of lithium-ion batteries".

Whittingham demonstrated the principle of electrochemical intercalation, in which lithium ions are stored inside the layered structure of an electrode without destroying it.

\subsection*{Significance and Impact}

His demonstration that lithium could be reversibly inserted into a host material established the rechargeable battery chemistry on which later lithium-ion batteries were built.

\subsection*{Later Career and Honors}

Whittingham became professor of chemistry and materials science at Binghamton University, where he continues research on energy storage.

\subsection*{Legacy}

Intercalation electrodes, introduced by Whittingham, are the basis of every rechargeable lithium-ion battery in use today.

\subsection*{Selected Publications}

"Electrical Energy Storage and Intercalation Chemistry" (Science, 1976)

\clearpage

\section{Akira Yoshino}

\subsection*{Early Life and Education}

Akira Yoshino was born in 1948 in Suita, Japan.

\subsection*{Career and Major Research}

Yoshino studied chemistry at Kyoto University, earning bachelor's and master's degrees in 1970 and 1972. He joined the Asahi Kasei Corporation, where he developed the carbon anode for lithium-ion batteries, using petroleum coke as the material that stores lithium.

\subsection*{Nobel Prize-winning Work}

The work for which Akira Yoshino was awarded the Nobel Prize in Chemistry in 2019 was described by the Nobel Committee as: "for the development of lithium-ion batteries".

Yoshino combined his carbon anode with Goodenough's cathode to build the first lithium-ion battery, and the technology was commercialized by Sony in 1991.

\subsection*{Significance and Impact}

His carbon anode removed the safety problems of metallic lithium and gave the battery its practical, rechargeable form, leading to the batteries used in consumer electronics.

\subsection*{Later Career and Honors}

Yoshino worked at Asahi Kasei for most of his career and later became a professor at Meijo University.

\subsection*{Legacy}

The lithium-ion battery he assembled now powers mobile phones, laptops, and electric vehicles across the world.

\subsection*{Selected Publications}

"The Birth of the Lithium-Ion Battery" (Angewandte Chemie International Edition, 2012)

\clearpage

\section*{Chapter Summary}

The 2019 prize recognized the development of the lithium-ion battery by three researchers: Whittingham established intercalation, Goodenough created the high-voltage cobalt oxide cathode, and Yoshino built the first practical cell with a carbon anode. Their rechargeable battery now powers portable electronics and electric vehicles worldwide.
